\chapter{Desenvolvimento}

Este capítulo detalha o funcionamento de algumas funcionalidades
desenvolvidas durante o projeto. Apenas algumas funcionalidades mais
chaves são abordadas neste capítulo, como o registro de humor, vinculo
de um gatilho, indicadores, e o registro de atividades.

\section{Processo de Login}

Logo ao abrir o aplicativo, o usuário é redirecionado para a tela de
login. A tela em si é relativamente simples, e contém os campos de
email e senha. Existe aqui um elemento proposital: uma forma de
garantir um certo nível de anonimato, em que o usuário não precisa se
identificar com um nome ou número de telefone, tampouco ativar a conta
através de uma plataforma de terceiro, como Google ou Facebook.

Isso se dá, em especial, por dois motivos. Primeiro, o uso de OAuth
exigiria uma conexão com plataformas de terceiro, que poderiam então
rastrear quais usuários têm acesso a qual aplicativo — algo que é
fartamente documentado na literatura sobre os mecanismos usados por
tais plataformas para rastrear e monitorar o acesso a dados de
usuários, com o fim de vender e monetizar essas informações.

Essa preocupação não é hipotética: um estudo que interceptou o tráfego de rede de
36 dos aplicativos mais populares de saúde mental (depressão e cessação de
tabagismo) encontrou transmissão de dados a pelo menos um terceiro em 92\% deles,
sendo que 81\% enviavam dados especificamente a serviços operados por Google ou
Facebook — e apenas 59\% desses casos declaravam esse compartilhamento em política
de privacidade \cite{Huckvale2019}.

Um estudo posterior, focado exclusivamente em 27 aplicativos de saúde
mental mais populares da Google Play Store, confirmou o mesmo padrão e
o caracterizou formalmente como uma ameaça de \textit{linkability}
(vinculabilidade) e \textit{identifiability} (identificabilidade): mesmo quando os
identificadores enviados a terceiros são pseudônimos (como \textit{advertising IDs}
ou hashes), eles ainda permitem que provedores de publicidade associem o
comportamento de um usuário através de diferentes aplicativos e serviços
\cite{Iwaya:22}. Como os próprios autores resumem, a mera vinculação de um
usuário a um aplicativo de saúde mental já pode revelar que essa pessoa possui
uma condição psicológica — seja ansiedade, depressão, ou, no caso de um
aplicativo como o Nexo, Transtorno Afetivo Bipolar —, tornando usuários desses
aplicativos particularmente vulneráveis a esse tipo de exposição
\cite{Iwaya:22}.

No caso do Nexo, essa cadeia de inferência é especialmente
sensível: a presença do aplicativo na lista de apps instalados de um usuário,
associada ao seu identificador publicitário, permitiria a uma plataforma como a
Meta inferir uma possível condição de saúde mental do usuário para fins de
segmentação de anúncios, sem que este tenha conhecimento ou consentimento
explícito disso.

Segundo, e retomando a proposta de oferecer ao usuário a possibilidade
de \textit{self-service}, aquele que decidisse manter uma instância
própria do aplicativo não precisaria de uma conta em uma plataforma de
terceiro, nem se cadastrar nelas.

\begin{figure}[H]
\centering
\caption{Tela de login}
\label{fig:screenshot-login}
\includegraphics[width=0.70\textwidth,height=0.60\textheight,keepaspectratio]{screenshot-login.jpeg}
\source{O autor, 2026}
\end{figure}

Por fim, destaca-se aqui que, conforme mencionado anteriormente na
modelagem de casos de uso (figura \ref{fig:use-cases}), ao realizar o
login pela primeira vez é enviado um email com um código de
verificação, seguindo a lógica dos \textit{magic links}. Isso adiciona
uma camada extra de segurança, garantindo que apenas o usuário que
solicitou o login possa acessar a conta.

\section{Página de Consulta de Registros}

A tela inicial é a tela nomeada apenas como ``Novo''. A proposta aqui
é ser a tela de cadastro rápido de um humor, e também é a tela que
visualmente mais destoa das demais. O usuário é apresentado com três
seções distintas:

\begin{itemize}
    \item Cadastro rápido de um humor
    \item Entre ontem e hoje
    \item Registros recentes
\end{itemize}

Respectivamente: uma seção em que o usuário é convidado a ``escolher
um humor'' e iniciar o processo de registro de uma nova entrada; a
seção em que são apresentados os indicadores de nível de energia
acumulado entre o dia anterior e o dia atual, além da média de sono do
dia comparada com a do dia anterior; e, por fim, a listagem dos dez
últimos registros de humor cadastrados pelo usuário.

\begin{figure}[H]
\centering
\caption{Tela de Consulta de Registros}
\label{fig:screenshot-consulta-registros}
\includegraphics[width=0.70\textwidth,height=0.60\textheight,keepaspectratio]{screenshot-initial-page.jpeg}
\source{O autor, 2026}
\end{figure}

Aqui já é introduzido ao usuário um padrão de organização e hierarquia:
todos os elementos centrais são organizados em \textit{cards}, exibidos
em no máximo duas colunas. Também é introduzido, caso o usuário já
possua registros de humor cadastrados, o padrão pelo qual esses
registros são resumidos em todas as seções do aplicativo: geralmente
um ícone, um título e alguma informação adicional em destaque.

\section{Criação de Registros de Humor}

Nesta seção, o usuário pode criar novos registros de humor, escolhendo
entre uma lista de opções pré-definidas. O humor selecionado na tela
anterior é mostrado aqui, e o usuário é convidado a julgar, em uma
escala de 1 a 10, como se sente em relação aos níveis de energia,
estresse e ansiedade, respectivamente.

\begin{figure}[H]
\centering
\caption{Criação de Registros de Humor}
\label{fig:screenshot-consulta-humor}
\includegraphics[width=0.70\textwidth,height=0.60\textheight,keepaspectratio]{screenshot-create-mood.jpeg}
\source{O autor, 2026}
\end{figure}

Naturalmente, é uma questão difícil de responder — e, como o próprio
aplicativo procura deixar claro, é importante que seja um momento
reflexivo, não necessariamente preciso.

O usuário pode adicionar uma anotação sobre aquele registro em
específico, padrão que se repete em outras telas do aplicativo.

Por fim, o usuário pode escolher o que compõe aquele registro, ou
seja, selecionar emoções e estados/sentimentos que sente naquele
momento. É apresentada uma listagem mista desses elementos, na qual o
usuário seleciona os que mais se aplicam a ele e atribui um nível de
intensidade para cada um. Aqui a escala é menor e mais simples, visando
facilitar a escolha e o julgamento do usuário, entre ``Leve'',
``Moderado'' e ``Intenso'', conforme a figura
\ref{fig:screenshot-add-components}.

\begin{figure}[H]
\centering
\caption{Adicionando sentimentos/emoções}
\label{fig:screenshot-add-components}
\includegraphics[width=0.70\textwidth,height=0.60\textheight,keepaspectratio]{screenshot-add-components.jpeg}
\source{O autor, 2026}
\end{figure}

Ao concluir esse processo e salvar, é apresentada ao usuário uma tela
de confirmação, na qual ele pode, opcionalmente, vincular o registro a
um gatilho específico — que pode ter ocorrido antes de ele salvar o
registro — ou criar um novo gatilho. O projeto novamente se ancora no
EMA, enriquecendo o registro de humor com mais contexto sobre o que
aconteceu.

\begin{figure}[H]
\centering
\caption{Vinculando um humor com um gatilho}
\label{fig:screenshot-link-mood-trigger}
\includegraphics[width=0.70\textwidth,height=0.60\textheight,keepaspectratio]{screenshot-link-mood-with-trigger.jpeg}
\source{O autor, 2026}
\end{figure}

Em seguida, o usuário pode indicar o impacto que aquele gatilho
exerceu sobre o registro de humor em questão. Novamente, trata-se de
uma estimativa subjetiva, sem pretensão de precisão.

\begin{figure}[H]
\centering
\caption{Avaliando o impacto de um gatilho no humor}
\label{fig:screenshot-assess-trigger-mood-impact}
\includegraphics[width=0.70\textwidth,height=0.60\textheight,keepaspectratio]{screenshot-link-mood-assess-impact.jpeg}
\source{O autor, 2026}
\end{figure}

Por fim, vale mencionar que, após esse processo, é enfileirado
internamente um \textit{background job} que recalcula, sob demanda, o
indicador de energia daquele dia, apresentado na figura
\ref{fig:screenshot-consulta-registros}. Diferentemente deste, os
indicadores semanais — como a correlação entre sono e energia,
descrita na seção \ref{sec:insights} — não são recalculados a cada
registro individual, sendo processados apenas pelo job semanal
agendado (seção \ref{sec:insights}).

\section{Indicadores}\label{sec:insights}

Nessa tela é apresentado um conjunto de indicadores construídos a partir
dos registros de humor, sono e gatilho do usuário. Diferente das demais
telas do aplicativo, que operam sobre dados em tempo real, os indicadores
são artefatos pré-computados: um \textit{worker} assíncrono (\textit{BullMQ})
percorre semanalmente todos os usuários ativos e enfileira, para cada
um, um conjunto de \textit{jobs} — um por tipo de indicador — que
processam a janela dos últimos sete dias e persistem o resultado em uma
tabela dedicada. Vale notar que nem todos os indicadores seguem esse
fluxo agendado: os indicadores diários (energia e sono) são recalculados
sob demanda, no momento em que o usuário registra um novo humor ou
sessão de sono, enquanto os indicadores semanais dependem
exclusivamente do \textit{job} de \textit{fan-out} de segunda-feira.

O aplicativo móvel apenas consulta esse resultado já
calculado, o que evita repetir cálculos custosos a cada abertura da tela
e desacopla a lógica analítica do restante da API.

\begin{figure}[H]
\centering
\caption{Tela de indicadores}
\label{fig:screenshot-insights}
\includegraphics[width=0.70\textwidth,height=0.60\textheight,keepaspectratio]{screenshot-insights.jpeg}
\source{O autor, 2026}
\end{figure}

% TODO: referenciar EMA / momentary assessment como justificativa da janela de 7 dias

São três os indicadores semanais exibidos: tendência de humor, correlação
entre sono e energia, e padrão de gatilhos.

\begin{itemize}
    \item \textbf{Tendência de humor}: compara a média de humor da
    primeira metade da semana com a da segunda metade, expressando a
    diferença como uma variação percentual. O resultado é classificado
    como melhora, piora ou estabilidade a partir de uma margem de $\pm10\%$.

    \item \textbf{Correlação entre sono e energia}: relaciona a duração
    do sono de uma noite com o nível de energia relatado na manhã
    seguinte, através do coeficiente de correlação de Pearson. Exige um
    mínimo de observações pareadas para ser calculado — caso contrário,
    o indicador simplesmente não é exibido.

    \item \textbf{Padrão de gatilhos}: identifica a categoria de gatilho
    mais frequente na semana e expressa sua predominância como percentual
    do total de registros.
\end{itemize}

É intencional, na figura \ref{fig:screenshot-insights}, que o bloco de
``Energia'' apareça vazio: o cálculo da correlação depende da existência
simultânea de registros de sono \emph{e} de humor em quantidade mínima,
e no período retratado o usuário ainda não havia acumulado dados
suficientes de ambos os tipos. Optou-se por não exibir uma correlação
calculada com poucos pontos, preferindo comunicar a ausência de dado a
apresentar um número pouco confiável.

O bloco maior no topo da tela, o ``Resumo Semanal'', condensa os
indicadores mais centrais — média de sono do período e um rótulo de
humor (``Estável'', no exemplo). Esse rótulo, no entanto, é o elemento
mais delicado de toda a funcionalidade. Embora exista uma fórmula capaz
de produzir uma classificação de estabilidade a partir da variância dos
registros de humor, atribuir esse tipo de afirmação a um usuário —
especialmente em um aplicativo que é utilizado por pessoas com
Transtorno Afetivo Bipolar, para quem uma leitura equivocada de
``estabilidade'' pode mascarar justamente o início de um episódio —
exige um embasamento mais cuidadoso do que uma regra estatística
simples. Por esse motivo, o rótulo apresentado nesta versão do
aplicativo é estático (não deriva de fato do cálculo de estabilidade) e
a funcionalidade pode ser inteiramente desabilitada pelo usuário nas
configurações. A decisão de manter ou não uma afirmação automatizada
desse tipo, e sob quais critérios, permanece um ponto em aberto que
mereceria revisão de literatura específica antes de ser tratado como
funcionalidade definitiva.

% TODO: literatura sobre riscos de auto-monitoramento automatizado de humor / mood-tracking apps e TAB

\section{Histórico}


Nessa tela é apresentado o histórico consolidado de registros de humor,
sono, gatilho e ações de cuidado (como sessões de terapia), organizados
cronologicamente e agrupados por dia, conforme a figura
\ref{fig:screenshot-history}. O usuário pode ainda filtrar a listagem por
categoria (``Humor'', ``Sono'', ``Gatilho'' e ``Ações de Cuidado''),
reduzindo o volume de informação exibida.

\begin{figure}[H]
\centering
\caption{Tela de Histórico}
\label{fig:screenshot-history}
\includegraphics[width=0.70\textwidth,height=0.60\textheight,keepaspectratio]{screenshot-history.jpeg}
\source{O autor, 2026}
\end{figure}

Como os quatro tipos de dado residem em recursos distintos da API, a
tela consulta os respectivos \textit{endpoints} em paralelo e, em
seguida, processa o resultado no próprio aplicativo. Para isso, é
utilizado um pequeno serviço que atua como \textit{mapper}, transformando
os dados de cada \textit{endpoint} em um formato comum, apresentado no
código \ref{lst:history-card}.

\newpage

\begin{lstlisting}[caption={Formato comum da tela de histórico}, label={lst:history-card}]
export type HistoryCard = {
  id: string; // "<category>-<id>" to avoid collisions
  category: HistoryCategory;
  subtype?: CareActionSubtype; // only set when category === "care_action"
  timestamp: string; // ISO string -- used for sorting and display
  title: string;
  summary: string;
  badge?: HistoryBadge;
  raw: MoodEntry | SleepRecord | Trigger | CareAction;
};
\end{lstlisting}

Após a normalização, os quatro conjuntos de dados são combinados em uma
única lista ordenada pelo campo \texttt{timestamp} e reagrupados por dia
para exibição, conforme visto na figura \ref{fig:screenshot-history}, nos
cabeçalhos ``terça-feira, 14 de julho'' e ``segunda-feira, 13 de julho''.
Essa unificação é proposital: ao invés de apresentar quatro telas
separadas por tipo de dado, o histórico busca reconstruir uma narrativa
única do dia do usuário, aproximando-se do princípio de \textit{Ecological
Momentary Assessment} já mencionado anteriormente — o contexto de um
registro de humor (por exemplo, o gatilho ``Interno'' registrado três
minutos antes, na mesma figura) fica visualmente próximo do evento que
possivelmente o motivou.

A propriedade \texttt{raw} armazena o dado original, permitindo que a
listagem redirecione o usuário para a tela de detalhe correta ao tocar em
um card, independentemente da categoria. Já a propriedade \texttt{badge}
reaproveita o mesmo mecanismo já descrito na seção de indicadores: para
um registro de sono, sinaliza se a duração foi insuficiente; para um
registro de humor, destaca a dimensão mais relevante daquele registro
(como ``Ansiedade alta'' na figura \ref{fig:screenshot-history}); para um
gatilho, indica sua categoria, e assim por diante.

\FloatBarrier
